\documentclass[12pt,a4paper]{article}

\usepackage[utf8]{inputenc}
\usepackage[T1]{fontenc}
\usepackage[spanish,es-tabla]{babel}
\usepackage{geometry}
\usepackage{enumitem}
\usepackage{array}
\usepackage{booktabs}
\usepackage{longtable}
\usepackage{hyperref}
\usepackage{microtype}
\usepackage{amssymb}

\geometry{left=2.5cm,right=2.5cm,top=2.8cm,bottom=2.8cm}

\hypersetup{
    colorlinks=false,
    pdfborder={0 0 1}
}

\setlist[itemize]{label=\(\blacksquare\), leftmargin=1.4em}
\setlist[enumerate]{leftmargin=1.6em}
\renewcommand{\arraystretch}{1.12}
\newcolumntype{L}[1]{>{\raggedright\arraybackslash}m{#1}}
\newcommand{\firma}{\rule{2.6cm}{0.4pt}}
\newcommand{\correo}[1]{\nolinkurl{#1}}

\title{Reglamento del grupo\\
\large Reglamento y Código de Ética\\
\normalsize ApuntesUCT, Taller de Integración, Semestre [I/II]-[AÑO]}
\author{}
\date{}

\begin{document}

\maketitle

\section{Introducción}

En el presente documento se establecen los principios, normas de convivencia, responsabilidades y sanciones que regirán al equipo del proyecto \textbf{ApuntesUCT}. Su propósito es asegurar una correcta coordinación entre los integrantes, mantener una ejecución ordenada del trabajo académico y reducir conflictos durante el desarrollo del semestre.

El reglamento aplica a los integrantes de los equipos INT2 e INT4, considerando que ambos grupos participan en un mismo proyecto y dependen de una coordinación constante para cumplir entregas, avances, integraciones y presentaciones.

\subsection{Integrantes del proyecto}

El proyecto está compuesto por 10 integrantes, separados en los equipos INT2 e INT4.

\subsubsection*{Equipo INT2}

{\small
\renewcommand{\arraystretch}{1.35}
\begin{longtable}{@{}L{0.05\textwidth}L{0.30\textwidth}L{0.29\textwidth}@{\hspace{0.45cm}}L{0.22\textwidth}@{}}
\toprule
\textbf{N} & \textbf{Nombre completo} & \textbf{Correo electrónico} & \textbf{Firma}\\
\midrule
1 & Gabriel Esteban Gutierrez Yañez & \correo{ggutierrez2024@alu.uct.cl} & \firma\\
2 & Ailyn Alejandra Melillán Bustos & \correo{amelillan2025@alu.uct.cl} & \firma\\
3 & Martina Alejandra Iturrieta Liempi & \correo{miturrieta2025@alu.uct.cl} & \firma\\
4 & Raúl Alejandro Rodríguez Flores & \correo{rrodriguez2025@alu.uct.cl} & \firma\\
5 & David Ignacio Villegas Araneda & \correo{dvillegas2025@alu.uct.cl} & \firma\\
6 & Antonio Ignacio Lara Fuentealba & \correo{alara2024@alu.uct.cl} & \firma\\
\bottomrule
\end{longtable}
}
\newpage
\subsubsection*{Equipo INT4}

{\small
\renewcommand{\arraystretch}{1.35}
\begin{longtable}{@{}L{0.05\textwidth}L{0.30\textwidth}L{0.29\textwidth}@{\hspace{0.45cm}}L{0.22\textwidth}@{}}
\toprule
\textbf{N} & \textbf{Nombre completo} & \textbf{Correo electrónico} & \textbf{Firma}\\
\midrule
1 & Marcelo Fabian Santana Astorga & \correo{marcelo.santana2023@alu.uct.cl} & \firma\\
2 & Eduardo Javier Escares Ortiz & \correo{eescares2023@alu.uct.cl} & \firma\\
3 & Yaninna Rossio Álvarez Álvarez & \correo{yalvarez2023@alu.uct.cl} & \firma\\
4 & Nelson Javier Quiñinao Isla & \correo{nquininao2023@alu.uct.cl} & \firma\\
\bottomrule
\end{longtable}
}

\section{Reglamento y convivencia}

En esta sección se detallan las normas de convivencia, comunicación y trabajo que el grupo asume para el correcto desarrollo del proyecto ApuntesUCT durante el semestre.

\subsection{Asistencia a reuniones}

\begin{enumerate}
    \item \textbf{Asistencia obligatoria:} Se espera que todos los integrantes asistan a las reuniones de trabajo acordadas por el grupo, ya sean presenciales o virtuales.

    \item \textbf{Notificación previa:} Las inasistencias deben avisarse con al menos 3 horas de anticipación a través de los canales oficiales definidos por el equipo. El aviso debe dirigirse al grupo o al Scrum Master correspondiente, indicando una razón válida.

    \item \textbf{Razón válida:} Se considerarán razones válidas aquellas asociadas a problemas de salud, emergencias familiares, dificultades de traslado justificadas, cruces académicos informados con anticipación u otras situaciones razonables evaluadas por el equipo.

    \item \textbf{Emergencias:} Si ocurre una emergencia que impida avisar con anticipación, el integrante deberá informar apenas le sea posible. En estos casos, la situación no se considerará automáticamente como falta, siempre que exista una explicación razonable.

    \item \textbf{Inasistencia sin aviso:} La ausencia a una reunión sin aviso previo, sin comunicación durante el evento y sin una justificación válida posterior podrá sumar una amonestación.

    \item \textbf{Atrasos:} Se espera puntualidad en reuniones, presentaciones y sesiones de trabajo. Los atrasos reiterados podrán considerarse falta leve.

    \item \textbf{Desconexión sin aviso:} En reuniones virtuales, desconectarse sin aviso y sin reincorporación posterior se tratará como una inasistencia no notificada, salvo que exista una causa justificada.
\end{enumerate}

\subsection{Comunicación y coordinación}

\begin{enumerate}
    \item \textbf{Canales oficiales:} El equipo deberá mantener canales claros para comunicación diaria, coordinación de tareas, avisos de inasistencia, bloqueos técnicos y decisiones relevantes del proyecto.

    \item \textbf{Reporte de avances:} Cada integrante debe comunicar oportunamente sus avances, dificultades y bloqueos al Scrum Master correspondiente o al equipo completo cuando la situación afecte a ambos grupos.

    \item \textbf{Transparencia ante problemas:} Si un integrante detecta que no podrá cumplir una tarea, entrega o plazo acordado, debe informarlo apenas lo identifique y no esperar hasta el momento de la entrega.

    \item \textbf{Responsabilidad compartida:} Si un integrante observa que otra persona no está trabajando o que el Scrum Master no está cumpliendo correctamente su rol de seguimiento, tendrá la responsabilidad de informar la situación al profesor si lo considera necesario.

    \item \textbf{Respeto y buena comunicación:} Toda comunicación debe mantenerse en términos respetuosos, tanto en reuniones presenciales como en canales digitales.
\end{enumerate}

\subsection{Reuniones entre INT2 e INT4}

\begin{enumerate}
    \item \textbf{Reunión semanal conjunta:} Cada lunes se realizará una reunión conjunta entre los equipos INT2 e INT4 para revisar avances, dependencias, bloqueos, decisiones técnicas y próximos compromisos.

    \item \textbf{Daily entre Scrum Masters:} Se considera posible realizar una daily breve entre los Scrum Masters de ambos equipos, con duración aproximada de 10 a 15 minutos, para mantener visibilidad del avance y detectar problemas de integración.

    \item \textbf{Reunión posterior a presentaciones de INT2:} Luego de cada presentación realizada por el equipo INT2, se realizará una reunión de coordinación entre ambos grupos. Esta reunión podrá ser solo entre Scrum Masters, sujeta a disponibilidad, salvo que exista una situación que amerite la participación completa de ambos equipos.

    \item \textbf{Disponibilidad del equipo:} La asistencia a reuniones extraordinarias dependerá de la urgencia del tema, la disponibilidad de los integrantes y el impacto que la situación tenga sobre el avance del proyecto.
\end{enumerate}

\subsection{Bitácora y seguimiento}

\begin{enumerate}
    \item \textbf{Responsabilidad del Scrum Master:} Es responsabilidad de cada Scrum Master mantener una bitácora de cada reunión realizada por su equipo o por ambos equipos.

    \item \textbf{Contenido mínimo de la bitácora:} La bitácora debe registrar fecha, asistentes, temas tratados, acuerdos, responsables, plazos y observaciones relevantes.

    \item \textbf{Registro de faltas:} Cada Scrum Master debe estar atento a posibles faltas de los integrantes de su equipo, registrando hechos de manera objetiva y con evidencia cuando corresponda.

    \item \textbf{Transparencia:} Las bitácoras y registros deben estar disponibles para el equipo cuando sea necesario revisar acuerdos, responsabilidades o situaciones de conflicto.
\end{enumerate}

\section{Protocolos de trabajo y entregas}

\subsection{Entregas de avance}

\begin{enumerate}
    \item \textbf{Frecuencia mínima:} Las entregas de avance deberán realizarse al menos 2 veces por semana, ya sea mediante commits, pull requests, actualizaciones documentadas, evidencias en la plataforma de trabajo o la herramienta acordada por el equipo.

    \item \textbf{Plazo para INT2:} Para el equipo INT2, el plazo final del segundo commit o avance semanal será el día lunes. Este margen permite disponer de al menos 1 día para corregir errores antes de la entrega o revisión correspondiente.

    \item \textbf{Plazo para INT4:} Para el equipo INT4, el plazo final del segundo commit o avance semanal será el día martes.

    \item \textbf{Evidencia de avance:} Todo avance relevante debe quedar respaldado mediante GitHub, documentos compartidos, bitácora, pull request, issue, tablero de tareas u otra herramienta acordada por el grupo.

    \item \textbf{No entrega de avances:} No entregar avances dentro de los plazos acordados, especialmente si afecta a otros integrantes o bloquea integración, podrá constituir falta leve, moderada o grave según el impacto y la reiteración.
\end{enumerate}

\subsection{Protocolos de trabajo en GitHub}

\begin{enumerate}
    \item \textbf{Uso de ramas:} No se debe trabajar directamente sobre la rama \texttt{main}. Cada integrante debe trabajar en ramas de desarrollo o ramas asociadas a tareas específicas.

    \item \textbf{Pull requests:} Los cambios relevantes deberán integrarse mediante pull requests cuando el flujo de trabajo del equipo lo contemple.

    \item \textbf{Revisión previa:} Antes de integrar cambios importantes, el equipo debe verificar que el código compile, funcione y no rompa el trabajo existente.

    \item \textbf{Push a main:} Los push o merge hacia \texttt{main} deberán ser realizados solamente por el Scrum Master correspondiente, y únicamente cuando el trabajo se encuentre funcional en la rama de producción o integración definida por el equipo.

    \item \textbf{Momento de integración:} La integración final hacia \texttt{main} se realizará en el día correspondiente al último plazo definido para cada equipo: lunes para INT2 y martes para INT4, salvo acuerdo explícito distinto.

    \item \textbf{Mensajes de commit:} Los commits deben tener mensajes claros y descriptivos del cambio realizado, evitando mensajes vacíos o poco informativos como \texttt{cambios}, \texttt{fix} o similares.

    \item \textbf{Integridad del código:} Está prohibido eliminar, sobrescribir o modificar el trabajo de otro integrante sin aviso previo y consenso del equipo.
\end{enumerate}

\section{Conducta general y trabajo en equipo}

\begin{enumerate}
    \item \textbf{Distribución equitativa de tareas:} Las responsabilidades deben repartirse de manera equilibrada, considerando carga de trabajo, habilidades, disponibilidad y dependencias entre INT2 e INT4.

    \item \textbf{Cumplimiento de responsabilidades:} Cada integrante debe cumplir las tareas que acepta o informar oportunamente si necesita apoyo, extensión de plazo o reasignación.

    \item \textbf{Colaboración entre equipos:} INT2 e INT4 deben coordinarse especialmente en contrato API, integración, pruebas, errores compartidos y cambios que afecten a Web, Mobile o backend.

    \item \textbf{Uso responsable de herramientas de IA:} En caso de utilizar herramientas de inteligencia artificial como apoyo, se debe respetar lo indicado por el profesor y comprender el contenido generado antes de incorporarlo al proyecto.

    \item \textbf{Plagio o copia sin comprensión:} Copiar código, documentación o soluciones sin comprender su funcionamiento, o sin dar crédito cuando corresponda, se considerará una falta grave.

    \item \textbf{Instancias evaluadas grupalmente:} En presentaciones, defensas u otras evaluaciones grupales, la ausencia debe justificarse mediante certificado médico u otro respaldo formal equivalente.
\end{enumerate}

\section{Propuesta de sanciones}

Se establece un sistema de sanciones simple, progresivo y orientado a resolver conflictos de manera transparente. Las faltas se clasifican en leves, moderadas y graves, según su impacto sobre el equipo, la entrega, la integración y la convivencia.

\begin{itemize}
    \item \textbf{Falta leve:} Atrasos reiterados, avisos tardíos de inasistencia, baja frecuencia de avances, mensajes de commit poco descriptivos, falta menor de comunicación o incumplimientos que no bloquean directamente al equipo. Se resuelve con llamado de atención interno y registro por parte del Scrum Master.

    \item \textbf{Falta moderada:} No informar avances o problemas relevantes, no cumplir protocolos de GitHub, no entregar avances comprometidos, faltar sin aviso ni justificación válida, reincidir en faltas leves o generar bloqueos evitables al equipo. Se registra con evidencia y puede notificarse al profesor si el equipo lo estima necesario.

    \item \textbf{Falta grave:} No asistencia reiterada sin aviso, no entrega de tareas que perjudique de forma importante al grupo, alterar o eliminar trabajo ajeno sin consenso, plagio, ausencia injustificada en instancias evaluadas, o reincidencia en faltas moderadas. Esta falta conlleva una instancia formal de resolución de conflicto.
\end{itemize}

\subsection{Escalamiento de faltas}

\begin{enumerate}
    \item Dos faltas leves equivalen a una falta moderada.
    \item Dos faltas moderadas equivalen a una falta grave.
    \item Una falta grave conlleva una resolución de conflicto, en la que el equipo, el integrante involucrado y el profesor podrán revisar antecedentes y definir medidas.
\end{enumerate}

\subsection{Resolución de conflicto}

Cuando una situación escale a falta grave o afecte directamente el avance del proyecto, se solicitará una instancia de conversación formal. En esta reunión se deberán revisar hechos, evidencias, bitácoras, historial de avances y responsabilidades asignadas.

Dependiendo de la evaluación del profesor y del equipo, la resolución podrá incluir redistribución de tareas, compromiso escrito, seguimiento especial, ajuste de responsabilidades, comunicación formal al profesor o las medidas académicas que correspondan según la asignatura.

\subsection{Consideraciones}

\begin{itemize}
    \item Las sanciones deben aplicarse con criterio, considerando contexto, evidencia, reiteración e impacto real sobre el proyecto.
    \item Las situaciones justificadas y comunicadas oportunamente no deben tratarse igual que una falta sin aviso.
    \item El objetivo principal del reglamento es prevenir conflictos y asegurar continuidad del trabajo, no castigar de manera arbitraria.
    \item Todo registro de faltas debe ser objetivo, verificable y conocido por el integrante involucrado.
\end{itemize}

\end{document}
